% Figura 2.1. Del parche al token en el Vision Transformer.
% Se usa en sec:marco_vision, subsección "Vision Transformer y Swin Transformer", después del
% párrafo que define el ViT. Cada parche de la fila superior y su token comparten color, para que se
% lea la correspondencia; el token de clase va en gris porque no viene de ningún parche.

\begin{figure}[H]
    \centering
    \begin{tikzpicture}[x=1mm, y=1mm]

        % Imagen dividida en una rejilla de 4x4 parches; la fila superior va coloreada.
        \draw[figtrazo, thick, fill=figrelleno] (-26, -12) rectangle (0, 12);
        \foreach \c [count=\i from 0] in {figazul, fignaranja, figverde, figpurpura} {
            \fill[\c!35] (-26 + \i*6.5, 6) rectangle (-19.5 + \i*6.5, 12);
        }
        \foreach \x in {-19.5, -13, -6.5} { \draw[figtrazo] (\x, -12) -- (\x, 12); }
        \foreach \y in {-6, 0, 6} { \draw[figtrazo] (-26, \y) -- (0, \y); }
        \draw[figtrazo, thick] (-26, -12) rectangle (0, 12);
        \node[font=\small, text=figtrazo, align=center] at (-13, -18) {Imagen dividida\\en parches};

        % Operación de proyección.
        \node[figtarea, text width=28mm] (proj) at (22, 0)
            {Proyectar cada parche\\a un token};

        % Token de clase y tokens de los parches, con su embedding de posición.
        \node[fignodo, minimum width=7mm, minimum height=6mm, inner sep=1pt, fill=figrelleno]
            (t0) at (50, 6) {\footnotesize[CLS]};
        \node[font=\footnotesize, text=figtrazo] at (50, -3) {$+$};
        \node[font=\footnotesize, text=figtrazo] at (50, -9) {$p_{0}$};
        \foreach \c [count=\i from 1] in {figazul, fignaranja, figverde, figpurpura} {
            \node[fignodo, draw=\c!70!black, fill=\c!35, minimum width=6mm, minimum height=6mm,
                  inner sep=1pt] (t\i) at (50 + \i*9, 6) {\footnotesize $x_{\i}$};
            \node[font=\footnotesize, text=figtrazo] at (50 + \i*9, -3) {$+$};
            \node[font=\footnotesize, text=figtrazo] at (50 + \i*9, -9) {$p_{\i}$};
        }
        \node[font=\small, text=figtrazo] at (95, 6) {$\cdots$};
        \node[font=\small, text=figtrazo, align=center] at (72, -18)
            {Secuencia de tokens con\\embedding de posición por parche};

        % Flechas.
        \draw[figflecha] (0, 0) -- (proj.west);
        \draw[figflecha] (proj.east) -- (43, 0) |- (t0.west);

    \end{tikzpicture}
    \caption{Del parche al token en el Vision Transformer, con el token de clase ([CLS]) al inicio.}
    \label{fig:marco_vit_patches}
\end{figure}
